% Add "answers" option to print the answers
\documentclass[addpoints]{exam}

\usepackage[utf8]{inputenc}
\usepackage{calc}
\usepackage[dvipsnames]{xcolor}
\usepackage{ifthen}
\usepackage{tcolorbox}

% \makebox[\textwidth]{Name:\enspace\hrulefill}
%
% \vspace{5mm}
%\makebox[\textwidth]{Instructor’s name:\enspace\hrulefill}

\definecolor{MyLightGray}{gray}{0.60}

\newcommand{\truefalse}
{		\begin{checkboxes}
			 \choice True
			 \choice False. Explanation:
		\end{checkboxes}
		\vspace{10pt}
}


% Check if 'answers' is an option
\newcommand{\ifanswers}[2]{%
  \ifthenelse{\boolean{exam@answers}}{#1}{#2}%
}



\newcommand{\answerspace}[1]
{
	\ifprintanswers
	\else
		\setlength{\mylength}{10pt}%
		% \vspace{#1\mylength}
		% Make the grey box
		\par\nobreak\noindent
		\makebox[\linewidth][l]{%
			% set inner padding to zero, so only the frame shows
			{\setlength{\fboxsep}{0pt}%
			% draw a color frame, white (or transparent) inside
			\fcolorbox{MyLightGray}{white}{%
			% the inner box has zero width content but fixed height
			\parbox[t][#1\mylength][c]{%
				% subtract twice the frame‐thickness from the full width
				\dimexpr\linewidth-2\fboxrule\relax
			}{\strut}%
			}%
			}%
		}%
		\par
		\vspace{3\mylength}

  	\fi
}
\newlength{\mylength}

\newcommand{\code}[1]{{\texttt{#1}}}


%\usepackage{listings}
\usepackage{color}
\usepackage{lastpage}
\usepackage{verbatim}
\usepackage{fancyvrb}
\usepackage{graphicx}
\usepackage{float}
\usepackage{minted}
\usepackage{multicol}

\newcommand{\includecode}[2][c]{\lstinputlisting[caption=, escapechar=, language=#1]{#2}}
\RecustomVerbatimCommand{\VerbatimInput}{VerbatimInput}{}

\renewcommand{\solutiontitle}{\noindent\textbf{Solution sketch: }}

\definecolor{SolutionColor}{rgb}{0.8,0.9,1}

\cfoot{Page \thepage\ of \pageref{LastPage}}

\rhead{Initials: \rule{20mm}{0.4pt}}



%%%%%%%%%%%%%%%%%%%%%%%%%%%%%%%%%%%%%%%%
\begin{document}
\fvset{numbers=left,numbersep=3pt}
\thispagestyle{empty}

\begin{center}
	{\LARGE CS 408 Fall 2026 -- Structural Coverage Quiz \\}
\fbox{\fbox{\parbox{5.5in}{\centering
	\begin{itemize}
        \item Try to come up with a solution on your own first, but
          after that talk to your friends and neighbors.

        \item Full points will be rewarded for \textbf{any} (reasonable)
          attempt.
        \end{itemize}
	}}}
\end{center}

\makebox[\textwidth]{\textbf{Name (please print)}: \rule{100mm}{0.8pt} }
%\end{center}

\vspace{10mm}
%\begin{center}
\makebox[\textwidth]{\textbf{PUID}: \rule{100mm}{0.8pt} \hspace{35mm}}
%\end{center}


\begin{questions}

\question Each program below contains a fault. For each one, name the
  \textbf{weakest} criterion we saw today (i.e., the lowest one in the
  subsumption chain) whose test requirements \emph{force} a test suite
  to execute the fault, and give a test case that reveals it. If no
  criterion does, say so and explain what you would use instead.

  \begin{parts}

  \part The specification is given in the comment.
  \begin{minted}{c}
// Effects: returns 1 if x > 0, and -1 otherwise
int sign(int x) {
  int s = 0;
  if (x > 0) {
    s = 1;
  }
  return s;
}
  \end{minted}
  \answerspace{8} %
  \begin{solution} %
    \textbf{Edge coverage.} The CFG has four nodes (\texttt{BB3},
    \texttt{C4}, \texttt{BB5}, \texttt{BB7}) and four edges. The single
    test \texttt{sign(5)} visits all four nodes and returns 1, which is
    correct -- so node coverage is satisfied by a suite that misses the
    fault entirely. Edge coverage additionally requires the false edge
    \texttt{(C4, BB7)}, which forces a test with $x \leq 0$; every such
    test returns 0 instead of $-1$. E.g.\ \texttt{sign(-3)}.

    This is the ``missing case'' shape: the bug is that \emph{nothing
    happened}, so there is no node to visit.
  \end{solution}

%   \part Assume the suite already contains \texttt{compute(1, -1)} and
%   \texttt{compute(-1, 1)}.
%   \begin{minted}{c}
% int compute(int x, int y) {
%   int z = 1;
%   if (x > 0) {
%     z = 0;
%   }
%   if (y > 0) {
%     z = x / z;
%   }
%   return z;
% }
%   \end{minted}
%   \answerspace{10} %
%   \begin{solution} %
%     \textbf{Edge-pair coverage.} The two given tests between them tour
%     every node and all seven edges, and neither divides by zero. Neither
%     decision is wrong on its own: the fault lives in the
%     \emph{interaction}, since the first branch establishes $z = 0$ and
%     the second then divides by it. Edge-pair coverage requires the
%     length-2 path \texttt{[BB4, C6, BB7]}, which can only be toured when
%     $x > 0$ and $y > 0$. E.g.\ \texttt{compute(1, 1)}, which throws.
%   \end{solution}

  \part The specification is given in the comment.
  \begin{minted}{c}
// Effects: returns the average (rounded down) of the elements of a.
//          Returns 0 if n == 0.
int average(int a[], int n) {
  int total = 0;
  for (int i = 0; i < n; i++) {
    total += a[i];
  }
  return total / n;
}
  \end{minted}
  \answerspace{8} %
  \begin{solution} %
    \textbf{Edge-pair coverage} -- not prime-path coverage, which is the
    tempting answer.

    Any single test with $n \geq 1$ tours every node and every edge,
    including the loop-exit edge, because a test that enters the loop
    must eventually leave it. So node and edge coverage are both
    satisfied without ever skipping the body.

    What forces the skip is the length-2 path
    \texttt{[FL1I, FL1C, BB7]} -- initialize, test, fall straight out --
    which is reachable only when $n = 0$. E.g.\ \texttt{average(\{\}, 0)},
    which divides by zero. Prime-path coverage also requires it (that
    path is prime), but it is not the \emph{weakest} criterion that does.
  \end{solution}

%   \part The specification is given in the comment.
%   \begin{minted}{c}
% // Effects: returns the average of a and b, rounded toward zero
% int midpoint(int a, int b) {
%   return (a + b) / 2;
% }
%   \end{minted}
%   \answerspace{10} %
%   \begin{solution} %
%     \textbf{No criterion helps.} The CFG is a single node with no edges,
%     so \emph{every} criterion we saw today -- up to and including
%     complete path coverage -- is satisfied by one arbitrary test case.
%     A suite can be 100\% covered under all of them, have a perfectly
%     good oracle, and still pass: \texttt{midpoint(2, 4)} returns 3 and
%     is correct.

%     The fault needs particular \emph{values}, not a particular path:
%     \texttt{a + b} overflows, so \texttt{midpoint(INT\_MAX, INT\_MAX)}
%     returns a negative number. Finding it is a job for input space
%     partitioning with boundary values on \texttt{a} and \texttt{b}.

%     The moral: structural criteria assume faults are visible in control
%     flow, and assume any two tests touring the same path are equivalent.
%     Both assumptions fail here.
%   \end{solution}
  \end{parts}

\end{questions}

\end{document}
